\chapter[变量的数学——从直观与思辨到成熟的数学科学]{变量的数学}
\begin{center}
{\large\sffamily\color{BellaGray}——从直观与思辨到成熟的数学科学}
\end{center}
\vspace{8mm}

读者都读过了一本通常的微积分教本，这样就会知道这是一门很有用的科学，尽管从这类教材中他很少能见到新的例子。再说，一门科学是否很有用也不是只靠几个例子能说明的。读者们会懂得了微积分中有许多解决问题的方法。如果不是遇到了很难的题目，或很细致的定理，微积分不是一门很难念的课程，而应该是很生动的。但是很多读者都对微积分的数学方法很不以为然。具体地说就是很不习惯 $\varepsilon$--$\delta$ 之类的语言，很不满意于对许多概念的过分仔细的分析。所以，本书打算这样开始：首先从历史发展的轨迹说明微积分何以有这样不“友好”的“界面”（users-unfriendly-interface）？但是本章并不是一个比较全面的微积分学历史的介绍，所以许多重要的人和事都没有讲。我们的目的只在于说明何以会有 $\varepsilon$--$\delta$ 这样令常人望而生畏的东西，说明这正是科学进步的结果，是数学科学区别于其它科学最明显的特点。本章结束以后，我们将再就微积分的若干主要领域介绍这种语言与方法是如何更有效地表述了微积分的主要思想，如何更有效地刻画了宇宙的规律，同时在这个过程中深化了自己，发展了自己（包括自己的语言与方法）。

“初等数学是常量的数学，高等数学是变量的数学”。这是老生常谈了，而且大体也是正确的。变量的数学在刻画自然界，乃至人类社会生活中取得了何等辉煌的胜利，这都是人所共知的了。但是什么是变量的数学？因为将“变”的概念引入数学又引起了何等深刻的变化？乃至于什么是“变”或“变量”？这些问题是值得我们去进一步思考的。我们将从历史的发展来看一下，这些问题是如何进入数学家的视野的。当代数学的一个最主要的起源地是希腊。希腊文明的所谓古典时期（即公元前6世纪至前3世纪），数学就已经形成了一个独立的学科。在那时，数学与哲学的关系是密不可分的，希腊人对许多数学问题的处理还有浓厚的思辨色彩。然而就是这样，关于变量和变化的数学问题已开始孕育了。简略地回溯一下这段历史，有助于我们去体会为什么微积分会有今天这样的形状，为什么我们不得不绞尽脑汁来对付 $\varepsilon$--$\delta$。也会体会到，两千多年前提出的问题至今仍未完全“解决”。但是，这样做，就不得不进入哲学的领域。这是我们力不能及的。所以，我们只能作一些浅显的介绍，而建议有兴趣的读者去读一些比较专门的著作。我们愿向读者推荐两本书：

\begin{quote}\small
M. Kline, \textit{Mathematics -- The Loss of Certainty}, Oxford University Press, 1980. 有中文译本：李宏魁译，《数学：确定性的丧失》，湖南科学技术出版社，1997。

T. Dantzig, \textit{Number -- The Language of Science}, George Allen \& Unwin Ltd, 1938. 有中文译本：苏仲湘译，《数：科学的语言》，上海教育出版社，2000。
\end{quote}

这两本书都是严肃的科普著作，不过要请读者注意不要以为读了它们就是读了哲学了。

希腊文明是唯一的这样一种古代文明：它承认人的理性的力量，人凭借着理性，再加上观察实验，就可以发现宇宙的规律，而不必求助于超自然的力量。希腊人较少受宗教的束缚，敢于反对传统，敢于反对教条的权威。对于他们，科学的任务就在于发现宇宙的规律。数学是科学的一部分（在希腊时期，几乎是科学的大部分），其任务也就是发现宇宙的规律。而且由希腊人，而伽利略，而牛顿……又总是认为宇宙的规律乃是数学的规律。这当然不是说，数学不为各个时代的经济发展的要求服务，不为技术服务，而是说，数学还有更深刻的任务，即探讨宇宙（原来主要是指自然界，但近几个世纪以来又越来越多地涉及人的社会生活）的规律。正因为它是科学，所以它又能更好地适应技术和经济发展的要求。对这个问题我们也不打算多作讨论，而只是指出，看到了这条主线就更容易理解微积分为什么会成为今天这样子。

数学作为一种演绎推理的科学，即数学中一切结论都应该用逻辑方法加以证明，按照罗素的说明，这个“原则”或者说是“规定”，是从毕达哥拉斯开始的。其后经由希腊时代的许多学派（主要是柏拉图学派）形成一个不可动摇的传统，而在欧几里得的《几何原本》中得到了完美的体现。毕达哥拉斯的数学又带有一种神秘的色彩，即他认为宇宙的本质是数（正整数，其实有理数也在其内，因为整数 $n$ 是以1为单位，重复 $n$——一个正整数——次而得，而有理数 $n/m$ 则是以 $1/m$ 为新单位重复 $n$ 次而得，这个新单位重复 $m$——又是一个正整数——次又可得原单位1）。既然如此，任何一个线段都应该由若干（正整数）个“单元”构成。这当然是原子论在数学中的反映。所以，任意两条线段 $l_1,l_2$ 都由相应的正整数 $m_1,m_2$ 个单元构成。单元是公共的单位，若两条线段 $l_1,l_2$ 有公共的单位，使它们分别由 $m_1,m_2$ 个单位构成（$m_1,m_2$ 是正整数），我们称这种情况为 $l_1,l_2$ “可公度”。但是其后毕达哥拉斯一位弟子正是利用了他所证明了的著名的毕达哥拉斯定理（中国称为勾股定理）发现，正方形的边长（设为1）与对角线长（用我们现在的记号是 $\sqrt2$）是“不可公度的”，即不可能找到任何的单元使1与 $\sqrt2$ 各含 $m$ 与 $n$（均为正整数）个单元，亦即不论 $m$ 与 $n$ 是什么样的正整数均不能使
\[
\frac{n}{m}=\sqrt2.
\]
用我们现在的语言来说，即 $\sqrt2$ 是无理数。这个证明应该是读者们所熟知的。按数学史家的研究，是欧几里得或其弟子们提出的。我们这里只想提到一点，即当时就已经用反证法来证明这个结论。

这是希腊数学中出现的一次危机。它的实质在于：有理数本质上只可用于刻画离散的对象，而几何图形如线段等等本质上却是连续的。但是什么叫连续，我们暂时还只能直觉地去掌握它。至少现在我们知道，用离散性的有理数不能刻画连续性的对象。希腊人还没有掌握无理数的理论，那是两千多年以后的事了。而希腊人对数学的要求又是很严格的，因此，可以说希腊人回避了数而专注于几何。他们研究量而回避数。例如线段之长是量，面积大小也是量。对于两个量例如线段 $l_1,l_2$ 之长可以讨论它们的“比”（ratio）。那么什么是比呢？在《几何原本》卷V中，给出了以下的“定义”：“比是两个同类量之间的一种大小关系”。今天看来，这实在算不上是一个“定义”，但《几何原本》中确实是这样说的。从现代数学的要求看来，《几何原本》在严格性方面是大有问题的，所以两个量之比是不是一个数是有问题的。但是紧接着讲到比例（即两个比之相等）原书上有一个定义5就很有意思了。因为照抄原书不太好懂，所以不妨用现代的语言说明如下：设有四个量，例如四条线段 $l_1,l_2,l_3,l_4$，如果对任意两个整数 $m$ 和 $n$，凡当
\[
ml_1\gtreqless nl_2
\]
时，必有
\[
ml_3\gtreqless nl_4,
\]
就说 $l_1$ 与 $l_2$ 之比等于 $l_3$ 与 $l_4$ 之比。这里要注意，$ml_1\gtreqless nl_2$ 并不是数的比较，而是长度的比较：“把 $l_1$ 首尾相连延长 $m$ 倍”大于“把 $l_2$ 首尾相连延长 $n$ 倍”，所以定义5是完全没有问题。如果我们承认量就是数，所以 $l_1$ 与 $l_2$ 之比就是 $l_1/l_2$，则这个定义讲的是：若由 $ml_1\gtreqless nl_2$ 必可得 $ml_3\gtreqless nl_4$，就说
\[
\frac{l_1}{l_2}=\frac{l_3}{l_4}.
\]
这个定义隐含地承认了“比”虽然不知是不是数，但是可以比较大小。如果我们把 $ml_1\gtreqless nl_2$ 定义为 $l_1$ 与 $l_2$ 之比 $\gtreqless n/m$（$n/m$ 确实是毕达哥拉斯也承认的数），则这个定义可以改述为：“若（$l_1$比$l_2$）$\gtreqless$（某数），则（$l_3$比$l_4$）也$\gtreqless$（同数），这时就说（$l_1$与$l_2$之比）=（$l_3$与$l_4$之比）”。

我们都知道，数（不论是否有理数）可以比较大小，希腊人则知道量可以比较大小。但是矛盾在于并不是所有的量都可以用数（实际上指有理数）来表示，例如单位正方形的对角线就不行。如果我们再多加一句话：能比较大小的就是数，则量就与数完全统一了。希腊人就是转不过这个弯子来，是不是希腊人笨？能创造如此辉煌的数学的民族会笨吗？转不过弯子来的症结何在？因为希腊人认为数毕竟要由“单元”构成，正如物体是原子构成的那样。原子是不可分的，所以单元就是不可分量；原子是一切整体的组成部分，所以单元作为一切整体的部分必小于一切数。这就是《几何原本》开宗明义的卷I公理5：“整体大于部分”。所以，这个“单元”就是“无穷小”。但是什么是无穷小呢？希腊人不仅是“不知道”，而且还可以证明“无穷小”不存在（下面我们要讲这件事）。所以希腊人无论如何也转不过这个弯子来。这个弯子一直到19世纪六七十年代才转过来。例如戴德金（J. W. R. Dedekind，1831--1916），用有理数的分割作为实数的定义，不妨说就是定义：凡能比较大小并作四则运算的东西就叫做实数（我们这样说还不尽符合戴德金的原意，详见本书第六章）。那么“无穷小”是什么呢？戴德金不需要无穷小，而是在实数理论基础上建立了极限理论，然后定义无穷小其实是极限为0的变量。这个弯子绕得太大了，不仅希腊人绕不过来，伽利略、牛顿都没有绕过来，而且从《几何原本》到戴德金，花了两千多年时间，所以如果说《几何原本》中的比例理论（一般认为应归功于欧多克萨斯（Eudoxus，公元前4世纪）\jiaokan{原文此处作“公元前3世纪”。欧多克萨斯约生于公元前408年、卒于公元前355年，应属公元前4世纪；且本章后文亦写其“生于公元前408年左右”。}）就是实数理论的前身，那又太有点“戏说”的味道。但是说从事后看来，希腊数学中已包含了其萌芽，恐怕是事实。

芝诺（Zeno of Elea，约公元前5世纪）和他的四个悖论更把关于数的“单元”（原子、不可分量、无穷小）所蕴含的深刻的矛盾突出出来了。芝诺其人生平不明，著作也没有流传下来。我们现在所见的来自亚里士多德的《物理学》。实际上，亚里士多德是企图解决芝诺所提出的矛盾的。关于芝诺的四个悖论说法很多，下面我们叙述的是罗素在《西方的智慧》（B. Russell, \textit{Wisdom of the West}, MacDonald \& Co., Ltd. 1959）一书中的说法。该书中文译者是马家驹、贺霖，世界知识出版社1992出版，48--50页。

芝诺的悖论是针对毕达哥拉斯的。前两个悖论针对的是一条直线是由无限多个单元构成的这样的论点。第一个悖论是阿基里斯（Achilles，希腊神话中的善跑者）永远追不上乌龟。阿基里斯的速度比乌龟快10倍，但是他让乌龟先走了100米。当阿基里斯追了这100米时，乌龟又向前走了10米；当阿基里斯追上这10米时，乌龟又向前走了1米。如此进行下去直至无穷，阿基里斯是永远追不上乌龟的。驳倒这一点并不难，因为若阿基里斯的速度为 $v$，则乌龟的速度是 $v/10$。当阿基里斯追到乌龟的出发点时，他用了时间 $100/v=T$。而在这个时间里，乌龟向前走了 $vT/10$。依此类推，阿基里斯所用的时间是
\[
T+\frac{T}{10}+\frac{T}{10^2}+\cdots
=T\sum_{n=0}^{\infty}\frac{1}{10^n}
=\frac{10}{9}T.
\]
这里涉及了无穷级数。很可惜，希腊人不懂得无穷级数，因此他们只能用思辨的方法来解释这些悖论。

第二个悖论是所谓“两分法”（dichotomy）。“阿基里斯”是用相对运动来反对直线是由无限多个单元组成这一论断的，“两分法”则说运动是不可能的，并以此反驳上述论断：若一个物体要沿某一直线走完一定的路程，它必先经过其中点；但在达到中点以前，先要到达 $1/4$ 点；在到达 $1/4$ 点以前先要到达 $1/8$ 点……仿此进行下去，该物体必须在有限时间内走完无限多个区段。而这是不可能的。因此，运动不可能。当然，用现代的数学知识来解释这个悖论也不困难。只要承认时间是无限可分的，则相应于空间的无限多个区段，也有时间的无限多个区段。而这些时间区段的长度成为一个几何级数，公比是 $1/2$，因而无限多个时间区段长之总和可以是有限的。这就解释了两分法悖论。

\begin{figure}[htbp]
\centering
\begin{tikzpicture}[bella figure,x=1.25cm,y=0.85cm,>=stealth]
\foreach \x/\lab in {0/A,1/B,2/C}{\fill (\x,2) circle (2pt) node[above=3pt] {$\lab$};}
\foreach \x in {0,1,2}{\fill (\x,1) circle (2pt);}
\draw[->] (2.2,1)--(2.65,1);
\foreach \x/\lab in {4/A,5/B,6/C}{\fill (\x,2) circle (2pt) node[above=3pt] {$\lab$};}
\foreach \x/\lab in {4.6/A',5.6/B',6.6/C'}{\fill (\x,1) circle (2pt) node[above=3pt] {$\lab$};}
\foreach \x/\lab in {3.4/A'',4.4/B'',5.4/C''}{\fill (\x,0) circle (2pt) node[above=3pt] {$\lab$};}
\draw[<-] (3.0,0)--(3.45,0);
\end{tikzpicture}
\caption*{图1-1}
\end{figure}

另两个悖论则是说：假设直线由无限多个单元组成，固然会引起如上的悖论，假设它是由有限多个单元组成，同样也会产生悖论。这里的悖论又是一组两个，第一个用相对运动，另一个则用一个物体的运动。第三个悖论称为“运动场”（stade）。假设空间的一个线段是有限多个单元组成的，一个时间区段则由有限多个瞬间（即时间的单元）组成。若有三排点 $A,B,C$ 如图1-1。第一排不动；第二排向右移动，由于单元总数为有限，所以点的移动距离只能是有限个单元，所用的时间区段也只包含有限多个瞬间。所以我们可以假设每个瞬间移动一个单元。第三排则向左移动一个单元。总之点的移动只能是在一个瞬间（即一个时间单元）内，移动一个点（即一个空间单元）。但是如图1，原来的 $A$ 点在一个瞬间以后，向右移动到 $A'$，向左移动到 $A''$，相距两个点的距离。于是 $A''$ 这一排点相对于 $A'$ 所经过的点数两倍于相对于 $A$ 所经过的点数，即一瞬间走了两点。但是按单元总数为有限的假设，一瞬间只能走一个点。由此可见，运动是不可能的。这一段话是直接从亚里士多德《物理学》一书中引用的。

第四个悖论是飞箭不动。飞行中的箭在每一个瞬间都“占有本身大小相等的位置一定的时间”。因此在这个瞬间，箭是不动的，而时间是由有限多个瞬时组成的。因此，在整个时间中，飞箭都是不动的。矛盾的产生在于假设了时间与空间都是由有限个单元组成的。用我们现在的观点来看，问题在于芝诺讲到了速度，如果用 $x(t)$ 表示物体的位置，则在两个不同瞬间 $t_1$ 与 $t_2$，有速度就是假设了
\[
v=\frac{x(t_2)-x(t_1)}{t_2-t_1}\ne0.
\]
因此 $x(t_2)-x(t_1)\ne0$，所以就必然有运动。芝诺说，因为飞箭在每一瞬间都有一定位置，这就意味着飞箭在这一瞬间是静止的，这个论断在逻辑上就没有根据。上面我们说的其实就是导数的概念。飞箭在每一瞬间都有确定的位置，即 $x(t)$ 有确定的值，这并不是静止，要静止，至少需要
\[
x'(t)=0.
\]
当然，希腊人没有这样的概念，芝诺的悖论有违于现代的科学，但是芝诺用这些悖论来反驳时间空间是由“单元”构成的，确实是击中要害的。

这些悖论表明，我们直觉中以为没有什么困难的概念如运动、连续、变化等等，如果深入地进行逻辑的分析，就会发现我们并不真正懂得它们。但是因为芝诺提出这些问题时完全是从思辨出发的，因此，对芝诺的反驳也必然是思辨的。其中最重要的是亚里士多德。他把无限分成两类：实无限（actual infinity）和潜无限（potential infinity），前者是指某一时刻“一下子”就掌握到无限，后者则是在一个时间的过程中才能逐步展开、实现的无限。例如一个线段可以无限地分割，但在任一个具体的时刻都不能无穷尽分割到头（因此不存在毕达哥拉斯所谓单元）；$1,2,3,4,\ldots$ 可以无限地数下去，时间可以不断流逝，这些都是潜无限。亚里士多德认为，芝诺的种种悖论都是针对的实无限。因此，亚里士多德认为如果只限承认潜无限，则不但不会有悖论，而且正是抓住了现实世界的本质，芝诺的困难也都迎刃而解。不仅如此，亚里士多德还指出一条直线的“连续性”就表现在：如果把它分成左右两段，则左段的终点正是右段的起点，二者融合为一；时间也是这样，0时0分0秒既是昨天的最终一瞬，也是今天的第一瞬，二者也融合为一。但是数就不是这样；以整数 $1,2,3,4,5,6,\ldots$ 而言，如果在4处分开，则4是左段的最末一数，5是右段的第一数，其间还差了许多。即令是有理数也是一样，$\sqrt2$ 正是 $1,1.4,1.41,1.414,\ldots$ 与 $2,1.5,1.42,1.415,\ldots$ 之间的空隙。数的集合究竟有没有连续性，一直没有解决。希腊人没有系统的无理数理论，这个理论出现在19世纪中后叶戴德金（Dedekind）与康托尔（Cantor）之手，正是回答了亚里士多德的问题。欧几里得一直只注重几何，而尽量回避数，原因大概在此。

芝诺的悖论在希腊哲学中影响极大，而后人对芝诺悖论的诠释、评论也众说纷纭，有人说甚至有过于对待《圣经》。以上我们只从数学的发展来看，就立刻看出，这里的关键在于“单元”究竟是什么？或者说，“无穷小”、“不可分量”究竟是什么？今天回顾历史，实在不能责怪希腊人回答不了这些问题。在世界上的一切古老文明中，只有希腊人如此深刻地提出了问题。一直到19世纪中后期的柯西（Cauchy）、魏尔斯特拉斯（Weierstrass）才给出了我们现在满意的回答。而在当时，至少是在柏拉图、欧几里得这里，数学家们持极为严谨的态度，实际上是“回避”无穷小，而他们所采用的方法，实际上是走了亚里士多德潜无限的路，而且实际上开辟了我们今天的极限理论的道路。

不少人认为希腊人这样的研究数学实在是脱离实际和烦琐哲学，而且播下了今天的数学之脱离实际的“罪恶”的种子，对此实难苟同。因为运动、连续、变化不是我们每天都看得见的吗？不是宇宙本性的一部分吗？因此不可避免地成为数学家研究的对象。何况，在希腊人看来，数和其它赋形的事物如日月星辰、山川草木的区别其实不大。只不过数学的方法与当时已经存在的科学：天文学、光学、静力学及至生理学等等不同，它不能依靠观察与实验，而只有逻辑推理，这当然又与毕达哥拉斯开创的传统相联系。看一看当时一些杰出的数学家的研究，就会发现他们实际上都是沿着亚里士多德潜无限的思想前进的。也许人们说这是亚里士多德的哲学思想的“指导作用”，但是我们没有任何证据来证明这种说法。宁可说，当时希腊人的认识水平就只达到了这一步。

我们要举出的杰出数学家之一是欧多克萨斯（Eudoxus）。他生于公元前408年左右，比亚里士多德早了好几十年，所以没有“接受”亚里士多德哲学“指导”的好福气。他是希腊古典时代最伟大的数学家之一，可是其生平与著作俱无流传。一般认为，欧几里得《几何原本》关于比例的理论，
关于穷竭法（早期的积分法）等等都是他的贡献。欧多克萨斯的工作只讲量（如长度、角度、面积、体积……）而不讲数，这是为了避免毕达哥拉斯关于无理数的困难。他应用了一个原理，后来人们认为应该看作是一个公理，因为阿基米德（Archimedes）也用过它，所以称为 Archimedes--Eudoxus 公理，见《几何原本》卷 X，命题 1（这里和以下凡引用《几何原本》处，均可在李文林主编《数学珍宝——历史文献精选》（科学出版社，1998 年出版）中找到）。用我们现代的语言来说，它是：设有两个量 $a>b>0$，从 $a$ 中除去其一半以上，再从其余除去其中一半以上，如此进行下去，经过充分多次以后，所余必小于 $b$。这个原理极为重要，因为它排斥无穷小量（不可分量）的存在；为简单起见，每次均取 $a$ 的 $\lambda$ 倍，$\frac12\le \lambda<1$，即 $\lambda a$，于是所余为 $(1-\lambda)a$，$n$ 次以后所余当为 $(1-\lambda)^n a$，这个命题说，只要 $n$ 充分大，必有
\[
(1-\lambda)^n a<b.
\]
其实，$\lambda$ 不必 $\geq \frac12$，$\lambda>0$ 即可。欧多克萨斯取 $\lambda\geq\frac12$ 是为了下面穷竭法所需。如果有一个量是无穷小量，即是小于一切量而又不是 $0$ 的量，则令上面的 $b$ 为无穷小量，而取 $a$ 为单位量 $1$，则一定有一个 $n$，使 $(1-\lambda)^n<b$，所以 $b$ 一定大于某个非 $0$ 的量 $(1-\lambda)^n$ 而不会小于一切非 $0$ 量。所以，无穷小量是不存在的。这一段证明显然有亚里士多德的潜无限的味道。

欧多克萨斯然后用它来求圆的面积，严格地说，他并没有求出圆的面积，而是证明了下面的

\textbf{定理}\quad 圆与圆之比等于其直径上的正方形之比。（《几何原本》，卷 XII，2。）

我们仍要比较仔细地用现代语言来叙述《几何原本》中的证明，因为这就是早期的积分学。设有两圆，其直径分别为 $d$ 与 $d'$，面积相应地为 $S$ 与 $S'$，则待证明的就是下面的比例式：
\[
S:S'=d^2:d'^2. \tag{1}
\]
或
\[
S/d^2=S'/d'^2. \tag{2}
\]
我们当然知道，这个公比是 $\pi/4$，但是注意到希腊人连 $\sqrt2$ 都不承认，又怎么承认 $\pi$ 呢？所以欧多克萨斯严格不讲这个公比是什么，但到了阿基米德则进一步计算了 $\pi$ 的近似值；他分别用圆的内接正 $96$ 边形和外切正 $96$ 边形算出 $\pi=3.14103$ 和 $\pi=3.14271$，因此，$\pi$ 的真值在二者之间。

这个定理的证明分为两部分，第一部分是证明圆可以内接正多边形去“穷竭”，即只要多边形的边数充分大，内接正多边形面积就与圆面积近似到任意要求的程度，因此，这个方法称为穷竭法。不过这个名词并非来自希腊人，而是来自 17 世纪的欧洲。证明如下，记圆面积为 $S$，内接正 $n$ 边形面积为 $P_n$，再作一个各边平行于此 $n$ 边形的外切正 $n$ 边形，记其面积为 $Q_n$，很明显
\[
P_n<S<Q_n. \tag{3}
\]
而且，只要 $n\geq4$，很容易看到 $P_n>\frac12Q_n$。《几何原本》中是用 $n=4$ 即用内接外切正方形来完成这个证明的。$S-P_n$ 由 $n$ 个相同的弓形构成。如果把内接正 $n$ 边形边数加一倍，即设 $AB$ 为其一边，取 $AB$ 之中分点 $C$，联结 $AC,BC$，用这两个边代替 $AB$ 即得正 $2n$ 边形，其面积为 $P_{2n}$。$S-P_{2n}$ 是 $2n$ 个相同的弓形。今证这 $2n$ 个弓形面积小于 $\frac12(S-P_n)$。为此，只看一个边 $AB$，并作相应的矩形 $ABED$ 如图 1-2，则 $\triangle ABC=\frac12\,\text{矩形 }ABED$，而弓形 $ABFCGA$ 的面积小于 $2\triangle ABC$。

\begin{figure}[htbp]
\centering
\begin{tikzpicture}[bella figure,scale=0.72]
\draw (0,0) circle (2.35);
\draw (-1.75,-1.55)--(1.75,-1.55)--(1.75,1.55)--(-1.75,1.55)--cycle;
\draw (-1.75,1.55)--(0,2.28)--(1.75,1.55);
\draw (-1.75,1.55)--(-1.75,2.28)--(1.75,2.28)--(1.75,1.55);
\draw (0,1.55)--(0,2.28);
\node[left] at (-1.75,1.55) {$A$};
\node[right] at (1.75,1.55) {$B$};
\node[above] at (0,2.28) {$C$};
\node[above left] at (-1.75,2.28) {$D$};
\node[above right] at (1.75,2.28) {$E$};
\node at (-0.95,1.92) {$G$};
\node at (0.95,1.92) {$F$};
\end{tikzpicture}
\caption*{图1-2}
\end{figure}

如果我们把内接正多边形边数增加一倍，则 $S-P_{2n}$ 与 $S-P_n$ 的关系是从（弓形 $ABFCGA$）中减去 $\triangle ABC$，而把原来的一个弓形换成两个小弓形，其弦各为 $AC$ 与 $BC$。但被减去的 $\triangle ABC>\frac12$（弓形 $ABFCGA$），所以这一个弓形被减去了“一大半”。每一个弓形如此，$n$ 个弓形之总和当然也如此，所以随着每一次把内接正多边形的边数增加一倍，$S-P_n$ 就会被减去“一大半”。仿此进行下去，利用卷 X 命题 1，即知，只要 $n$ 取得足够大，$S-P_n$（其实就是 $|S-P_n|$）就会变得“要多小有多小”。

至此，读者们一定会说，由此确知，对任意 $\varepsilon>0$ 必定有 $N$ 存在，使当 $n>N$ 时，$|S-P_n|<\varepsilon$，亦即 $\lim_{n\to\infty}P_n=S$。但是不要太急，不要忘记希腊人还没有极限的明确概念。而他们对数学严格性的要求又使他们不能用没有定义过的东西。不过，他们实际上并没有真正做到这一点，例如什么是圆面积，他们就只说它是一个量，而不是数；但是什么是量，什么是量的比都没有下定义。在 20 世纪初希尔伯特又写了一本《几何基础》（科学出版社有中译本），这样才达到了 20 世纪数学严格性的要求，但是已把《几何原本》改得面目全非了。不使用极限，使希腊人绕了一个大弯——求助于反证法。

证明的第二部分是求证（1）式，这里他们用了反证法。设（1）不成立，则一定可以找到比例第四项 $S''$ 使
\[
S:S''=d^2:d'^2. \tag{4}
\]
关于比例第四项是否一定存在，《几何原本》又是默认而未证明的。如果 $S''<S'$，则取 $n$ 充分大，如上所述可以使 $S'$ 的内接正 $n$ 边形 $P_n'$ 与 $S'$ 之差任意小，从而小于 $S'-S''$，即
\[
|P_n'-S'|<|S'-S''|,
\]
于是有
\[
S'>P_n'>S''. \tag{5}
\]
在《几何原本》中紧紧放在这个定理前还有一个命题：圆内接相似多边形之比等于圆直径平方之比，这又是很容易证明的。故
\[
P_n:P_n'=d^2:d'^2. \tag{6}
\]
再由（4），有
\[
P_n:P_n'=S:S''.
\]
因 $P_n<S$，故 $P_n'<S''$，但这与（5）矛盾。如果设 $S''>S'$，也同样会出现矛盾。定理证毕。

我们不厌其烦地引证了二千多年前的文献，是为了与我们现在都懂得的微积分作一比较。我们甚至不妨说，19 世纪中叶，柯西和魏尔斯特拉斯正是复兴了欧多克萨斯的穷竭法，才完成了微积分基础的奠定。首先，他们正式给出了极限的定义，这是欧多克萨斯没有做到的；有了极限，他们就以 $P_n$ 的极限作为圆面积的定义，这也是欧多克萨斯没有想到的；有了这一切，他们就再也用不着绕反证法的大弯子了。他们定义无穷小即是极限为 $0$ 的变量，从而“真正”地排除了毕达哥拉斯“单元”以及后来的不可分量的“痕迹”。这只是在欧多克萨斯的基础上走了更明确的一步。他们也宣告了亚里士多德潜无限思想的“胜利”，把笼罩着的思辨色彩“清除”了。这一点恐怕他们自己也没有想到。千真万确的是，他们把变量、极限等等归结为不等式，确是得了欧多克萨斯的真传。可惜的是，正是这一点成了 $\varepsilon$--$\delta$ 的罪名，因为它们使人望而生畏了，这却是变量的数学的精华。

不过，凡事有利必有弊。《几何原本》在数学上的严格性成了人们的“模范”以后，也必然会对人的直觉、感觉、不严格的推理有所排斥。而这种排斥时常也把许多极为生动的、通向真理的其它道路封锁了。以亚里士多德而言，潜无限的思想被推到了主导的地位，实无限的思想难道就一无是处了吗？否，历史证明不是这样。《几何原本》一个最明显的缺点是轻视数学的实际应用，这不是说希腊数学整个都是轻视应用的。即以欧多克萨斯而言，他不仅是一位数学家，还是一个天文学家（他曾经提出一个关于天体运动的模型），他还是一个医生、一个地理学家。另一位更著名的大人物是阿基米德（Archimedes，287--212，B.C.）。他是亚历山大时期（或称希腊化时期）的希腊数学和科学的代表人物。他较晚于欧几里得，但二人大体并不重叠；他不仅在数学，而且在力学、机械学、光学等各方面都有重大贡献。他还是一位伟大的发明家，而欧几里得主要贡献则是完成了《几何原本》（其中不少是他的弟子们的功绩），有点像孔夫子说的“述而不作”，而阿基米德的创造天才却是光彩夺目的。因为我们的目的只是讨论微积分的发展与方法的演变，所以现在就不来讨论希腊数学各个时期特点的比较及其功过了。好在到了公元前 2 世纪至前 1 世纪，罗马人征服了希腊，出现了罗马帝国，希腊数学也退到历史的后台去了。

千年沉睡以后，当数学科学再次醒来，已经开始进入资本主义时代了，社会经济的发展要求变了。到了 16 和 17 世纪，人类面临新的数学问题。其中一些重大问题是：求一般几何形体的长度、面积、体积；求曲线的切线；求运动的速度与加速度；求函数的极大与极小等等。我们的读者一眼就可以看出来，当时人类需要的新科学就是微积分。这些问题一方面有明确的实用背景，但又常与哲学上的论争，与人类争取从天主教神学以及亚里士多德的教条下获得思想解放的斗争紧密联系。伽利略在比萨斜塔上的实验（但是科学史的研究表明，他并没有做这个实验，而是用斜面的实验反驳了亚里士多德），围绕日心说的斗争是大家都熟知的了。新时代呼唤巨人，他们是笛卡儿、哥白尼、伽利略、开普勒直到牛顿和莱布尼茨。下面是一些例子。

第一个例子是关于积分学的。那个时代人们努力的重点是发展希腊人的穷竭法，但是他们又丢了穷竭法中的精华：用复杂的反证法绕过缺少极限理论带来的困难，把问题归结于不等式，从而为 19 世纪的 $\varepsilon$--$\delta$ 铺平了道路。他们反而又回到了不可分量（亦即毕达哥拉斯的单元），不过更大胆，而且几乎没有多少玄妙的思辨色彩了。开普勒（Johannes Kepler，1571--1630）有一本书名叫《测量酒桶体积的新科学》，顾名思义，其中会讲到面积和体积。在开普勒看来，圆是无穷多个顶角无穷小的扇形堆成的，球是无穷多个互相平行的厚度无穷小的圆柱形堆成的。伽利略也是这样，例如他认为若一物体（应为质点）作等加速运动，则在 $t=0$ 到 $t=T$ 这段时间走过的距离将由 $\triangle OAB$ 之面积来表示。因为，若在时刻 $t$，速度为 $A'B'$，伽利略将它乘以无穷小的时间间隔，$\triangle OAB$ 就是由无穷多个直线段 $A'B'$ 组成的，因此他得出结论：距离由 $\triangle OAB$ 之面积来表示（图 1-3）。

\begin{figure}[htbp]
\centering
\begin{tikzpicture}[bella figure,x=1cm,y=0.75cm,>=stealth]
\draw[->] (0,0)--(6.4,0) node[right] {$t$};
\draw (0,0)--(5.4,4.1);
\draw (5.4,0)--(5.4,4.1);
\draw (2.9,0)--(2.9,2.2);
\node[below] at (2.9,0) {$t$};
\node[below] at (5.4,0) {$T$};
\node[left] at (0,0) {$O$};
\node[right] at (5.4,0) {$A$};
\node[right] at (5.4,4.1) {$B$};
\node[above left] at (2.9,2.2) {$B'$};
\node[below] at (2.9,0) {$A'$};
\end{tikzpicture}
\caption*{图1-3}
\end{figure}

这样看来，伽利略和开普勒在面积体积问题上是一样的，采用了不可分量的观点。这种观点在伽利略的学生卡瓦列里（Bonaventura Cavalieri，1598--1647）的著作里表现得最系统。简单地说，他认为线段是由点组成的，正如链子是由珠子穿成的一样；面是由线组成的，正如布是由纱织成的一样；立体是由平面组成的，正如书是一页一页的纸叠起来的一样。点、线、面就是相应的不可分量，它们在一个或几个维数上是无穷小。

所以，如果说欧多克萨斯是在千方百计地回避无穷小，而为此费尽心思利用了 Archimedes--Eudoxus 公理，利用了不等式，利用了反证法，这一个时代的数学家却是硬着头皮往前闯。虽说不严格，却得到了许多正确的结果，这样不会招致人们的批评吗？卡瓦列里的回答是绝妙的。他说，我的目的只是为了改善穷竭法；我的方法确实不严格，但是我的结果很有用，有用就行，严格不严格那是哲学家的事，别的几何学家不是和我一样不严格吗？

首先，它确实改善了穷竭法。上面我们讲到的希腊的穷竭法太依赖于特殊的几何曲线（例如圆）的特殊性质，现在却有了一定的程式。例如要计算一段抛物线 $y=x^2$ 下的面积 $OAB$ 时（图 1-4），我们把 $OA$ 分成 $n$ 等分，每一段的长记为 $d$，于是相应的纵坐标是 $d^2,(2d)^2,(3d)^2,\ldots$。如果把 $n$ 取为无穷大，那么图上用实线画的区域都成了直线段，我们用虚线把它们补成矩形的。如果 $n$ 是无穷大，这些矩形和实线所成图形是没有区别的，它们就是不可分量。现在把这些不可分量的面积加起来，就得到 $OAB$ 的面积：
\[
\begin{aligned}
OAB
&= d(d)^2+d(2d)^2+\cdots+d(nd)^2\big|_{n=\infty}\\
&=d^3[1^2+2^2+\cdots+n^2]\big|_{n=\infty}\\
&=\left(\frac{OA}{n}\right)^3[1^2+2^2+\cdots+n^2]\big|_{n=\infty}\\
&=OA^3\left(\frac{2n^3+3n^2+n}{6n^3}\right)\bigg|_{n=\infty}.
\end{aligned}
\]
这里只用了一个很初等的代数公式
\[
1^2+2^2+\cdots+n^2=\frac16 n(n+1)(2n+1).
\]
还有 $\left.\frac1n\right|_{n=\infty}=0$。前式并不难证，后面的 $\left.\frac1n\right|_{n=\infty}=0$ 按卡瓦列里的说法，交给哲学家去管好了。不管怎样，许多人都可以按这个程式去计算积分了。如果把严格性暂置一旁，这确实是很大的进步。

\begin{figure}[htbp]
\centering
\begin{tikzpicture}[bella figure,x=0.9cm,y=0.6cm,>=stealth]
\draw[->] (0,0)--(7,0) node[right] {$x$};
\draw[domain=0:2.5,smooth,variable=\x] plot ({2.3*\x},{0.72*\x*\x});
\foreach \k in {1,...,6}{
  \pgfmathsetmacro{\xx}{0.9*\k}
  \pgfmathsetmacro{\yy}{0.72*(\xx/2.3)*(\xx/2.3)}
  \draw[dashed] (\xx-0.9,0) rectangle (\xx,\yy);
}
\node[left] at (0,0) {$O$};
\node[below] at (5.4,0) {$A$};
\node[right] at (5.75,4.5) {$B$};
\node[below] at (0.9,0) {$d$};
\node[below] at (1.8,0) {$d$};
\end{tikzpicture}
\caption*{图1-4}
\end{figure}

第二个例子是关于速度。速度和加速度这些运动学概念的研究在当时的重要性我们已在上面讲过。如果运动轨迹用 $x=x(t)$ 来表示，则大家知道 $\Delta x/\Delta t$ 表示平均速度。现在要求瞬时速度，可是又不能令 $\Delta t=0$，因为那样一来就连运动也没有了（牛顿语），所以牛顿称 $\Delta x/\Delta t$ 为最初比，然后让 $\Delta t$ 逐渐趋近 $0$（请注意，牛顿在这里不但使用了趋近的说法，还多次用了“极限”一词，只不过那时并没有明确的极限概念，牛顿也没有花功夫去说明极限的含意是什么），于是最初比有一个“极限” $\dot x(t)$，牛顿称为“最终比”；牛顿把变化中的量如 $x(t)$ 称为流量（fluent），而把其变化率 $\dot x(t)$ 称为流数（fluxion）。牛顿的方法论似乎有些混乱，有时他也使用不可分量，但是例如在《原理》一书中他又批评不可分量，明确地说：“如果我说某量由粒子组成，……不要以为我是指不可分量，而是指趋于 $0$ 的可分量”。牛顿以为用运动学的观点去解释这些问题就可以避免古希腊人的思辨方法、形而上学方法的困难。他做到了吗？下面再引牛顿的几段话（这些话都见于《原理》中文本第 38 页）：

“可能有人反对，认为不存在将趋于 $0$ 的量的最后比值，因为在量消失之前，比率就不是最后的，而当它们消失时，比率也没有了”（就是说，$\Delta t\ne0$ 时就不是瞬时速度，$\Delta t=0$ 就连运动也没有了）。“最后速度……就是……物体到达其最后处所并终止运动时的速度……最后比可以理解为……它消失的那一瞬间的比……在运动尚存的最后时刻速度有一极限，不能超越，这就是最后速度。”牛顿还说“我论及最小的、将消失的或最后的量，读者不要以为是在指确定大的量，而是指作无止境减小的量”。所谓瞬间是什么？“它消失那一瞬间”是什么？什么叫“无止境减小的量”，是不是 $0$？可不可以略去不计？这才是问题所在。

看一下切线问题还会更清楚。下面我们讲一下费马的方法，图 1-5 画出了求 $y=f(x)$ 之图像在 $P$ 点的切线的方法，此图含意自明。下面除了其中关键的 $PR$ 记为 $E$ 之外均不再解释。令 $G$（即 $x$ 点）得一增量 $E$（就是我们常用的 $\Delta x$）而变为 $G_1$，由于 $\triangle TGP\sim\triangle PRT_1$，所以
\[
TG:PG=GG_1:T_1R=E:T_1R.
\]
但是，当 $E$ 非常小时，费马说 $T_1R$ 和 $P_1R$ 相差不多，从而可以把 $T_1P_1$ 略去不计，用我们熟悉的语言就是当 $\Delta x$ 很小时，$dy$ 与 $\Delta y$ 之差可以忽略不计，所以
\[
TG:PG=E:(P_1G_1-PG)=E:[f(x+E)-f(x)].
\]
但 $PG=f(x)$，故上式给出
\[
TG=\frac{Ef(x)}{f(x+E)-f(x)}
=f(x)\bigg/\frac{f(x+E)-f(x)}{E}. \tag{7}
\]
最后费马说，“去掉 $E$ 项”（即令 $\Delta x=0$），立即可得 $TG$，在确定了 $T$ 点位置以后，连接 $TP$ 即得切线。

\begin{figure}[htbp]
\centering
\begin{tikzpicture}[bella figure,x=1cm,y=0.85cm,>=stealth]
\draw[->] (0,0)--(6.2,0);
\draw (0.4,0)--(4.1,4.2);
\draw[domain=2.7:5.1,smooth,variable=\x] plot ({\x},{0.20*(\x-1.8)^2+1.6});
\coordinate (P) at (2.7,2.6);
\coordinate (G) at (2.7,0);
\coordinate (G1) at (4.6,0);
\coordinate (P1) at (4.6,3.2);
\draw (P)--(G);
\draw (P)--(P1);
\draw (P1)--(G1);
\draw (4.1,4.2)--(4.6,3.2);
\node[left] at (0.4,0) {$T$};
\node[left] at (P) {$P$};
\node[below] at (G) {$G$};
\node[below] at (G1) {$G_1$};
\node[right] at (P1) {$P_1$};
\node[right] at (4.1,4.2) {$T_1$};
\node at (3.8,2.6) {$E$};
\node[right] at (4.6,2.6) {$R$};
\end{tikzpicture}
\caption*{图1-5}
\end{figure}

这里遇到的问题是：令 $E=0$ 时，（7）之分母成了 $0/0$ 而没有意义。然而要注意，在那个时代，人们知道的函数不多，例如笛卡儿就只讨论 $f(x)$ 是多项式的情况。这时 $f(x+E)-f(x)$ 仍是多项式，而且每一项都有因子 $E$，因此，（7）之分母中 $E$ 可以约去（注意，这当 $E=0$ 时是不合法的），因而不会出现 $0/0$ 问题了。但是令 $E=0$ 不但是做了不合法的事，而且 $E=0$ 就没戏可唱了，这和上面讲的 $\Delta t=0$ 就没有运动如出一辙。$E$ 又是 $0$ 又不是 $0$，是一个方生未死似 $0$ 而又非 $0$ 的无穷小量！可是什么是无穷小呢？再说上面略去了 $T_1P_1$，究竟什么样的东西可以略去不计呢？所以，情况仍与芝诺时代差不多，必须要回答什么是无穷小，而且不能只用思辨的方法，要解决具体问题，如当年用穷竭法计算圆面积一样。

我们不能像卡瓦列里说的那样，把严格性完全交给哲学家。《几何原本》对这个时代的数学家仍然是不可改变的规范。例如牛顿的基本著作《自然哲学的数学原理》就完全是模仿《几何原本》的。牛顿知道其中的问题时常出在无穷小量上面；说它不是零，并用看起来很像零的英文字母 $o$ 来记它，又常用 $a+o=a$，或在两个无穷小乘积中令它为零；说它是零，又允许以它为分母，并且称 $0/0$ 为不定式。18 世纪的许多卓越的数学家都为此绞尽了脑汁。这当然瞒不过哲学家的法眼。这里必须提到爱尔兰克罗因地方的主教乔治·贝克莱（Bishop George Berkeley，1685--1753）的一本著名的小册子《分析学家，或致一位不信神的数学家》（\textit{The Analyst, or a Discourse Addressed to an Infidel Mathematician}），读者可以在前面引用的李文林主编的《数学珍宝》一书 316--321 页上看到此文的摘译，下面引用的译文均来自此书。“不信神的数学家”是指哈雷（E. Halley），哈雷彗星的发现者，对牛顿写出《原理》一书起了重要的作用，并资助出版。伯克莱写这本书当然是为了宣扬自己的宗教观和自己的哲学。大体说来，他的论据就是：连牛顿的微积分、无穷小量那样模糊不清、逻辑混乱的东西都可以相信，为什么你们却不肯相信上帝呢？我们不想讨论哲学，但是应该指出，伯克莱的批评是切中要害的。伯克莱直接指名责备牛顿“曲线求积术”一文中给 $x$ 以非零的增量 $h$，然后作了一些代数运算，然后又叫 $h=0$，这岂非背离了逻辑学中最起码的矛盾律（即 $A$ 不是非 $A$）吗？而在神学中是不允许这样做的。而且他两次引用牛顿在该文中的原话：“在数学中最微小的误差也不可以忽略”来讽刺牛顿。伯克莱指出，牛顿的“流数”（即我们今天讲的导数）是不可理解的，更何况二阶、三阶……流数“全部超出整个人类理解能力”，这种又是零又不是零的增量，还有增量的增量……“这就是我们的现代数学家们……全部思想的基石”。可是这些全是不可理解的东西，至于由此得到的成功，伯克莱说：“他使用流数，就像建筑中使用脚手架一样，房子盖好了自然就立即将它们弃之一边。”“我所非议的不是您的结论，而是您的逻辑和方法；您是怎样进行证明的？您所熟知的对象是什么？关于它们您的表述是否清楚？您依据的原理是什么？它们是否可靠？您是如何应用它们的？”（着重点是本书作者加的——作者）至于得到了正确的结论，这是“因为这个错误被另一个相反的但程度相当的错误抵消了。”（这句话是指莱布尼茨。）伯克莱正确地点出了要害：那个似零又不是零的 $h$ 是什么？“这些消逝的增量又是什么呢？它们既不是有限量，也不是无限量，又不是零，难道我们不能称它们为消逝量的鬼魂吗？”其实，从我们上面引述的牛顿在《原理》一书中对自己的最终比的说明，并且一再强调自己的方法与不可分量并不相同。看来，难道不能认为牛顿自己也感到了类似的问题，因而理不直、气不壮，处于彷徨与无奈之中吗？

把微积分放在一个严格的基础上是很长一段时间数学家着重努力的大问题，但是有一点应该指出，这一段时期所有的数学家是同时用极大的力量从事扩大与深化微积分的应用的。现在的大学数学系，一直到高年级的许多课程的内容都是在 18 世纪、19 世纪形成的，这一点与希腊时代非常不一样。这是生产力迅猛发展及其带来的科学的迅猛发展的表现。达朗贝尔有一句名言，很可以反映当时的情况：“前进吧，你会得到信心！”

确实是这样，数学进入了一个高歌猛进的时代。新思想、新方法层出不穷，新领域一天被发现、被征服。然而却全处在伯克莱大主教提出的一长串问题的笼罩下。我们想从欧拉的《无穷小分析引论》\footnote{欧拉这一著作，是用拉丁文写的，书名 \textit{Introductio In Analysis Infinitorum}，1748。应译为《无限量分析引论》，1988 年，J. D. Blantan 有英译本，就采用了这样的书名：\textit{Introduction to Analysis of the Infinite}，我们仍然从俗把欧拉这部著作译为《无穷小分析引论》}中转述一大段，而且仍然沿用欧拉的记号。这一段正是讲的指数函数与对数。

任取一数 $a$ 为底，则 $a^0=1$。若 $\omega$ 表示一个无穷小，则 $a^\omega=1+\Psi$，$\Psi$ 也是一个无穷小。所以我们可设 $\Psi=k\omega$，而
\[
\omega=\log_a(1+k\omega). \tag{8}
\]
例如当 $a=10$ 时，$k=2.30258$。对任意数 $j$ 有
\[
a^{j\omega}=(1+k\omega)^j
=1+\frac{j}{1}k\omega+\frac{j(j-1)}{1\cdot2}k^2\omega^2+\cdots.
\]
令 $j=z/\omega$ 而 $z$ 为定数，则 $j$ 等于无穷大，代入上式有
\[
a^z=\left(1+\frac{kz}{j}\right)^j
=1+\frac{1}{1}kz+\frac{j(j-1)}{1\cdot2}\frac{k^2z^2}{j^2}
+\frac{j(j-1)(j-2)}{1\cdot2\cdot3}\left(\frac{kz}{j}\right)^3+\cdots. \tag{9}
\]
但因 $j$ 是无穷大，所以 $\frac{j-1}{j}=1,\frac{j-2}{j}=1$ 等等，代入上式，
\[
a^z=1+\frac{kz}{1}+\frac{k^2z^2}{1\cdot2}+\cdots. \tag{10}
\]
令 $z=1$ 有
\[
a=1+\frac{k}{1}+\frac{k^2}{1\cdot2}+\cdots. \tag{11}
\]
上面说到 $\Psi=k\omega$，取 $a$ 为相应于 $k=1$ 的那个数，则由（11）式有
\[
a=2.71828182845904523536028. \tag{12}
\]
（请注意，欧拉时代没有计算机，他怎么能用手算出 $e$ 到 23 位小数？）我们记这个数为 $e$。下面一段是欧拉的原话：

“为简单计，我们用符号 $e$ 表示此数：$e=2.718281828459\cdots$。它是自然对数或称双曲对数的底，它相应于 $k=1$，而且 $e$ 表示无穷级数 $1+\frac11+\frac1{1\cdot2}+\frac1{1\cdot2\cdot3}+\frac1{1\cdot2\cdot3\cdot4}+\cdots$ 之和”。

再看对数，由（9），令其中的 $a=e$，从而 $k=1$，有
\[
e^z=\left(1+\frac{z}{j}\right)^j
=1+\frac{z}{1}+\frac{z^2}{1\cdot2}\frac{j}{j}\frac{j-1}{j}+\cdots. \tag{13}
\]
如果把 $e^z$ 写成 $1+x$，则由上式
\[
\begin{aligned}
z&=j\big[(1+x)^{1/j}-1\big]\\
&=j\left[\frac{x}{j}+\frac{x^2}{1\cdot2}\frac1j\left(\frac1j-1\right)
+\frac{x^3}{1\cdot2\cdot3}\frac1j\left(\frac1j-1\right)\left(\frac1j-2\right)+\cdots\right]\\
&=x-\frac{x^2}{2}+\frac{x^3}{3}-\cdots.
\end{aligned}
\]
但是 $z$ 就是 $\log_e e^z=\log_e(1+x)$，所以又有
\[
\log_e(1+x)=x-\frac{x^2}{2}+\frac{x^3}{3}-\cdots. \tag{14}
\]

我们能同意伯克莱大主教说的：“所有这些错误碰巧都抵消了，所以欧拉得出了如此惊天动地的结果”吗？当然不能。欧拉是一位极为罕见的大数学家，他在处理数学与力学问题上的能力是非凡的。如果数学中以欧拉命名的东西都没有了，数学史一定要改写。但是数学仍然会发展，而且一定会和今天一样好。欧拉的成就应该归功于他非凡的洞察力和数学运算能力。更深一层来看，牛顿、莱布尼茨真正伟大之处在于他们真正抓住了变量的数学的本质，形成了高屋建瓴势如破竹的局面。所以在一两百年之内巨人辈出也正是时势造英雄，余下的工作就是把问题都搞清楚，否则数学是不可能再进步了。请读者看欧拉的一段文章目的就是请大家注意，欧拉没有一个结论是错的，但是真把他的毛病都消除，却是要大费周章。特别是由二项级数，再逐项求极限决非轻而易举。读者们读过的微积分教本上面那一小段引文要花多少篇幅才说得清楚，读者自己都有体会。

上面我们完全没有讲莱布尼茨的贡献。这不是由于在“发明”微积分的“知识产权”问题上我们有什么偏爱，而是由于，一开始我们就说了，本章并不是一个初步但比较全面的微积分历史。莱布尼茨和牛顿虽然在微积分的理论和方法上都不一样，但是在什么是微分 $dx,dy$（这都是莱布尼茨的记号）与不可分量、无穷小量的关系上，他与牛顿遇到相同的困难。伯克莱大主教在反对他们二人上似乎也是一视同仁的，所以我们就完全略去这一部分。好在在本章开始时介绍给读者的两本书中都有了介绍。

解决遗留下来的问题——其核心从希腊时代起就是关于无穷小量与不可分量等等的问题——出路在于极限理论。极限的概念牛顿等许多人都已经有了，而在柯西以前，最接近当代极限概念的是达朗贝尔。然而所有这些论述都过于几何化：所谓一个变量 $x$ 趋于 $a$，人们自觉或不自觉地是从运动学的角度来领会的。因此柯西最大的功绩在于把极限的研究算术化了。这方面的先驱应该指出的还有波尔查诺（Bernhard Bolzano，1781--1848）和阿贝尔（Niels Henrik Abel，1802--1829）。前者曾指出高斯关于代数的基本定理（每个代数方程都有根存在）证明不够严格，因为高斯利用了连续函数的中间值定理未加证明。他又是第一个给出不可求导的连续函数的例子（1848）的数学家，比魏尔斯特拉斯早 30 年。他多年来是布拉格的一位教士而未为人知。阿贝尔很年轻就去世了，一生贫病交加。他对当时分析中缺少严格性可说是痛心疾首，真正成就和影响最大的当然是柯西。他写了几本教本：《工科大学分析教程》（\textit{Cours d'Analyse de l'École Polytechnique}, 1821）、《无穷小计算概要》（\textit{Résumé des Leçons sur les Calculs Infinitésimals}, 1823）与《微分学讲义》（\textit{Leçons sur les Calculs Différentiels}, 1829），使我们看到了今天我们熟知的微积分教本。柯西所作的事就是使极限摆脱几何的束缚。极限概念如果不能摆脱几何的束缚，终将会引起混乱。例如说圆内接正 $n$ 边形当 $n\to\infty$ 时极限就是圆，那么就可以问，是否终于变成了圆？圆的性质全蕴含在正多边形的性质之中了？柯西明确指出，极限就是一个数。所谓变量 $x$ 趋向 $a$ 时 $f(x)$ 以 $A$ 为极限就是：若代表变量 $x$ 的一串数值无限地趋近固定数 $a$ 且其差可以任意地小时，$f(x)$ 与 $A$ 之差也可以任意小。但是稍后的数学家如魏尔斯特拉斯仍然对此定义不满，认为仍然包含了几何学与运动学的痕迹，应该更进一步算术化。变量是什么？魏尔斯特拉斯给的回答就是：若有一个无限非空集合 $D$（例如实数集），而 $x$ 可以取 $D$ 中任一元为值，$x$ 就称为一个变量。$\lim_{x\to a}f(x)=A$ 的意思就是：任给 $\varepsilon>0$ 必可找到 $\delta=\delta(\varepsilon)>0$，使当 $|x-a|<\delta$ 时，$|f(x)-A|<\varepsilon$。我们看到，这样做其实正是发展了《几何原本》中把穷竭法归结为不等式的实践，于是许多问题自然化解了：这些正 $n$ 边形最后会不会与圆重合？根本不必问这样的问题，反正有 $\lim_{n\to\infty}P_n=S$。每一个 $P_n$ 均不是 $S$，$|P_n|$ 中没有最后一个元。“最终比”是不是运动停止的那一瞬间的速度？无所谓这一瞬间那一瞬间，既无所谓运动开始的最初比，也无所谓运动停止“那一瞬间”的最终比。我们看到的只有 $\Delta x/\Delta t$ 是 $\Delta t$ 的函数，$\Delta t$ 与 $\Delta x/\Delta t$ 都是变量，而当 $\Delta t\to0$ 时 $\Delta x/\Delta t$ 也有极限 $\lim_{\Delta t\to0}\Delta x/\Delta t$，它就叫做瞬时速度。运动可以还继续，无所谓最终。什么是无穷小？它没有任何特殊的神秘之处，用不着与“单元”、“原子”、“不可分量”……扯到一起，它只是以 $0$ 为极限的变量而已。

先是柯西，再继之以魏尔斯特拉斯，把整个微分积分放在极限的基础上，而后来又放在不等式的基础上，可以说，魏尔斯特拉斯如愿以偿地“清洗”了运动学的“痕迹”，变量的数学静态化了（既已静态化，也就用不着再为芝诺的悖论操心了）。这就是读者们学过了的微积分。

于是我们看到，什么是导数？是极限；什么是积分？也是极限；什么是 $\sum a_n$？还是极限。（只有什么是微分？这比较麻烦，第三章全章就是为了回答这个问题。）什么是极限？也很清楚，极限是数。那么什么是数？这就追溯到古希腊的第一次危机。请看康托尔（G. Cantor）的回答：“$\sqrt3$ 只是一个尚未确定的数的符号，而不是它的定义。然而，用我的方法可以满意地给出其定义为
\[
\{1.7,1.73,1.732,\ldots\}.
\]
康托尔的话意思是，实数需要定义。其方法之一是定义实数为一个具有特定性质的有理数序列。这样一来，单位正方形的对角线长也是一个“数”，其定义是一个序列：
\[
\{1.4,1.41,1.414,\ldots\}.
\]
用这样的方法解决了毕达哥拉斯的困惑，使得直线上每一个点都对应一个数，每一个数也对应一个点，即使得点与点的坐标一一对应。这是解析几何的基础，而且把整个几何“算术化”了。例如两个线段的比现在很容易定义了，若 $l_1,l_2$ 的长是有理数，则 $l_1:l_2=l_1/l_2$ 也是有理数；若 $l_1,l_2$ 之长有一个是无理数，$l_1:l_2=l_1/l_2$ 仍有意义，无非可能是无理数而已。总之，量与数的对立消除了。不仅如此，有关极限理论最困难的定理也都得到了证明。这就是魏尔斯特拉斯所企盼的把整个微积分算术化。我们将在第六章详细讨论这些问题。

总之，微积分现在已经建立在稳固的基础上了。它有如一个“平台”，在这个平台上可以建立许多大厦，可以用它来解决许多以前所不能设想的问题。我们这本书就想对这个平台上的一些“大建筑”作一番漫游。当然，数学科学还有许多其它的平台，即具有自己的新思想、新方法的领域。我们虽然不一定能去参观，不一定能去“定居”，但是可以告诉读者，这些平台中大多数与牛顿的平台有密切的关系，而且从数学的历史来看，如果想找一个比牛顿平台更基本的，可能就只有《几何原本》了。

不过还有一个问题要说几句，既然微积分已建立在可靠的基础上了，是否由此就完全否定了直觉的作用？我们不打算进入哲学的领域。从数学的教学和研究来看，我们还只能说，这么多世纪的数学发展证明了人们的直觉还是可靠的，需要的只是深入的分析。直觉仍是数学的思想与理论的来源之一。例如芝诺的悖论，我们还只能说，如果不是终日冥想，而是进入具体的数学与物理学的领域，人们就不会只作这种抽象的玄想。从芝诺悖论会产生集合论的问题，也都在集合论作为数学的一个分支的框架内，或者得到解决，或者得到深化。又如无穷小、运动等等，我们只要知道了在数学中是如何处理这些问题的，就会发现这些直觉不但无害而且有益。而且，数学懂得越多，从直觉里得到益处也越多。至少是受“严格性不足”之害越少。所以直到现在，人们还是说“设某个增量 $\Delta x$ 是无穷小”，说什么某个小体积是“无穷小元”，或者“微元”之类的话，本书也是这样做的。总之，越往后走，越深入数学和物理学，您会得到越大的“自由”，而不必谨小慎微。

大约二百年前，达朗贝尔用一句鼓舞人的口号来描述当时的数学：“前进吧，你就会有信心”。而现在的情况则是：人类正充满信心，人类正在前进！
